% ======================================================================
% SECTION 10: ROBOTIC EXECUTION AND PHYSICAL AI GATEWAY
% File: sections/robotic_execution.tex
% ======================================================================

[PROCESSING INSTRUCTION: Generate this section (approximately 3-4 pages) covering the robotic
execution gateway that makes this a Physical AI sponsor system. Use the following source files:

SOURCE FILES TO PROCESS:
- national-platform/new\_paper/final\_paper/sections/psl\_usl\_standards.tex (PSL and USL frameworks)
- national-platform/new\_paper/final\_paper/sections/patient\_journey.tex (robotic procedures)
- patient-journey/paper/patient\_journey\_paper.tex (da Vinci Xi lobectomy, Franka Emika automation)
- national-platform/new\_paper/final\_paper/sections/site\_establishment.tex (site robot infrastructure)
- national-platform/new\_paper/final\_paper/sections/introduction.tex (Physical AI definition)

SUBSECTION STRUCTURE:

\subsection{Robot Execution Gateway (robot\_execution\_gateway)}
- Central interface between sponsor agents and physical robotic systems
- Task-order generation from sponsor decision to robot instruction
- Site capability matching: match procedure requirements to available robots
- Safety gate evaluation before every physical procedure
- Procedure-state management: pre-op, active, post-op, recovery
- Emergency stop handling and failover protocols

\subsection{Robotic Capability Registry}
- Registry of all qualified robots across trial sites
- USL scoring for each robot (9 robots evaluated in national platform)
- PSL compliance scoring for site-level readiness
- Capability matching algorithm: procedure to robot qualification
- Maintenance and recertification tracking

\subsection{Surgical Robot Operations}
- da Vinci Xi integration: surgical procedure orchestration
  (lobectomy example from patient journey, USL 87.5)
- Pre-operative digital twin simulation validation
- Intraoperative telemetry monitoring and adaptive control
- Post-operative outcome documentation

\subsection{Collaborative Robot Operations}
- Franka Emika integration: laboratory automation, pharmacy dosing (USL 88.75)
- Sample handling and processing workflows
- Dose preparation and dispensing verification
- Continuous calibration and quality assurance

\subsection{Device Telemetry and Procedure Provenance}
- Real-time telemetry ingestion from all robotic systems
- Procedure timeline reconstruction from sensor data
- Immutable procedure records for regulatory submission
- Integration with trust layer for provenance chains

PYTHON SCRIPT REQUIREMENTS (to be generated by Claude Code Opus 4.6 in the future):
- File: sponsor/template/scripts/robot\_execution\_gateway.py
  Purpose: Central robotic execution gateway. Implement task-order generation,
  site capability matching, safety gate evaluation (multi-step approval
  with configurable timeout), procedure-state machine (pre-op, active, post-op,
  recovery), and emergency stop handling. This is a critical Physical AI component.
- File: sponsor/template/scripts/robotic\_capability\_registry.py
  Purpose: Robot registry with USL scoring. Track all qualified robots, their
  capabilities, maintenance status, and USL scores. Implement capability
  matching: given a procedure type, return ranked list of qualified robots.
  Reference the 9 robots evaluated in the national platform paper.
- File: sponsor/template/scripts/surgical\_robot\_interface.py
  Purpose: da Vinci Xi surgical robot interface. Implement procedure orchestration
  (pre-op planning, intraoperative monitoring, post-op documentation), telemetry
  ingestion (force/torque, instrument position, insufflation pressure), and
  digital twin validation interface. Reference patient journey stage 5.
- File: sponsor/template/scripts/cobot\_interface.py
  Purpose: Franka Emika collaborative robot interface. Implement laboratory
  automation workflows (sample handling, centrifugation, aliquoting), pharmacy
  dosing (weight-based calculation, preparation verification, dispensing),
  and calibration management. Reference patient journey stages 7 and 8.
- File: sponsor/template/scripts/procedure\_provenance\_tracker.py
  Purpose: Immutable procedure provenance. Record every robotic action with
  timestamps, sensor readings, operator approvals, and outcome data. Generate
  regulatory-grade procedure reports. Interface with trust layer blockchain.
- All scripts must pass ruff lint (line-length 120, E/F/W rules).
- Process with Claude Code Opus 4.6 (1M token context).

TABLE REQUIREMENTS:
- Table 17: Robotic System Registry (9 robots from national platform)
  Columns: Robot | Manufacturer | Type | USL Score | Trial Function | Safety Level
  Include: da Vinci Xi, Hugo RAS, Versius, Franka Emika, Kinova Gen3,
  xArm7, Atlas, Optimus, Digit
- Table 18: Procedure Safety Gate Protocol
  Columns: Gate | Trigger | Required Approvals | Timeout | Fallback Action

FORMATTING NOTES:
- This section is central to the Physical AI concept. It must clearly establish
  that the sponsor system can both make decisions AND issue bounded, safety-checked
  instructions to physical robotic systems.
- Cross-reference \S\ref{sec:safety} for robotic safety monitoring.
- Cross-reference \S\ref{sec:trust} for provenance architecture.]

\subsection{Robot Execution Gateway}

[PROCESSING INSTRUCTION: Generate 1-1.5 pages on the robot\_execution\_gateway. This is the
critical component that makes this a Physical AI sponsor (not just a software sponsor).
Cover: task-order generation (translating sponsor agent decisions into bounded robotic
instructions), site capability matching (procedure requirements mapped to available
qualified robots), safety gate evaluation (multi-step approval with mandatory human
gates for high-risk procedures), procedure-state management (finite state machine:
idle to pre-op to active to post-op to recovery to complete), and emergency stop handling
(immediate halt, state preservation, failover to manual, incident documentation).
Process with Claude Code Opus 4.6 (1M token context).]

\subsection{Robotic Capability Registry}

[PROCESSING INSTRUCTION: Generate 1 page with Table 17 listing all 9 robots evaluated
in the national platform. For each robot, provide: manufacturer, type (surgical/cobot/humanoid),
USL score, primary trial function, and safety classification. Include the capability
matching algorithm description. Reference psl\_usl\_standards.tex from the national
platform paper. Process with Claude Code Opus 4.6 (1M token context).]

\subsection{Surgical Robot Operations}

[PROCESSING INSTRUCTION: Generate 0.75 page on da Vinci Xi integration. Reference the
168-minute robotic lobectomy from the patient journey (R0 resection, negative margins).
Cover pre-operative digital twin simulation, intraoperative telemetry, and post-operative
documentation. Process with Claude Code Opus 4.6 (1M token context).]

\subsection{Collaborative Robot Operations}

[PROCESSING INSTRUCTION: Generate 0.75 page on Franka Emika integration. Reference the
laboratory automation and pharmacy dosing from the patient journey (USL 88.75).
Cover sample handling, dose preparation, and continuous calibration.
Process with Claude Code Opus 4.6 (1M token context).]

\subsection{Device Telemetry and Procedure Provenance}

[PROCESSING INSTRUCTION: Generate 0.5 page on telemetry and provenance. Cover real-time
data ingestion, procedure timeline reconstruction, and immutable records.
Cross-reference the trust layer (\S\ref{sec:trust}).
Process with Claude Code Opus 4.6 (1M token context).]
